\documentclass[12pt,a4paper,spanish]{book} %%%Esto indica el tipo de documento.
%Va a ser un libro (book), el tama�o es a4, la lengua castellano (spanish)%%%
\usepackage{babel} %%%Incluimos el paquete Babel
%que sirve para separar correctamente
%las palabras de multitud de idiomas%%%
\usepackage[latin1]{inputenc}
%%%Este paquete permite poner acentos directamente%%%
\usepackage{amsmath}%%%Macros AMS%%%
\usepackage{amsthm}%%%Macros AMS para teoremas%%%
\usepackage{amsfonts}%%%Permite usar fuentes AMS%%%
\usepackage[dvips]{epsfig} %%%Inclusi�n de figuras postscript% con visualizaci�n posterior%%%\usepackage{indentfirst}%%%Espaciado de
\usepackage{titlesec}
\usepackage{lettrine}
\usepackage{eso-pic} 
\usepackage{xcolor}

\let\oldcleardoublepage\cleardoublepage
\renewcommand{\cleardoublepage}{\newpage{\pagestyle{empty}\oldcleardoublepage}}

%Cabecera%
\date{}
\usepackage{fancyhdr}
\setlength{\headheight}{15pt}
\pagestyle{fancyplain}
\newcommand{\diezpuntos}{\fontsize{8pt}{\baselineskip}\selectfont}
\addtolength{\headheight}{1\baselineskip}
\lhead{\diezpuntos ITISistemas}
%\chead{\diezpuntos }%
\rhead[\diezpuntos \bfseries \leftmark]{\diezpuntos \bfseries \rightmark}
\renewcommand{\headrulewidth}{0.5 pt}
\renewcommand{\footrulewidth}{0pt}
\lfoot{\diezpuntos}


%Comandos%

%*Marquesina*%
\newcommand\BackgroundPic{ 
\put(-5,0){ 
\parbox[b][\paperheight]{3cm}{% 
\vfill 
\includegraphics[width=3cm,height=\paperheight, 
keepaspectratio]{img/CircuitoMarquesina.eps}% 
\vfill 
}}} 

%Listas sin enumerar%
\renewcommand{\labelitemi}{-}
\renewcommand{\labelitemii}{-}
\renewcommand{\labelitemiii}{-}
%Listas enumeradas%
\renewcommand{\labelenumi}{%
\textbf{\theenumi}.$\star$
}
\renewcommand{\theenumii}{\arabic{enumii}}
\renewcommand{\labelenumii}{%
\textbf{\theenumi}.\theenumii.$\star$
}
\renewcommand{\theenumiii}{\arabic{enumiii}}
\renewcommand{\labelenumiii}{%
\textbf{\theenumi}.\theenumii.\theenumiii.$\star$
}

%Citas%
\let\oldquote\quote
\renewcommand\quote{\par\singlespacing\small\oldquote}
\let\oldquotation\quotation
\renewcommand\quotation[1]{\oldquotation\hfill\small\emph{\color{blue! 10!black}{#1}}}
\let\oldverse\verse
\renewcommand\verse{\par\singlespacing\small\oldverse}

%Capital%
\renewcommand{\LettrineFontHook}{\fontfamily{ptm}\mdseries\upshape\color{blue! 30!black}}

%Imagenes%
\newcommand{\imgIncl}[3]{
  \begin{figure}
		\centering
		\includegraphics[width=#1cm,angle=#2]{#3}
		\label{fig:#3}
	\end{figure}
	
	}
	
	%Parte%
\newcommand{\bigrule}{\titlerule[0.5mm]}
\titleformat{\part}[display] % cambiamos el formato de los cap�tulos
{\bfseries\Huge} % por defecto se usar�n caracteres de tama�o \Huge en negrita
{% contenido de la etiqueta
\titlerule % l�nea horizontal
\filleft % texto alineado a la derecha
\Large{Parte}  % "Parte" en tama�o \Large en lugar de \Huge
\Large\thepart} % n�mero de cap�tulo en tama�o \Large
{0mm} % espacio m�nimo entre etiqueta y cuerpo
{\filleft} % texto del cuerpo alineado a la derecha
[\vspace{0.5mm} \bigrule \vspace{4cm}] % despu�s del cuerpo, dejar espacio vertical y trazar l�nea horizontal gruesa

%Capitulo%
\titleformat{\chapter}[hang] % cambiamos el formato de los cap�tulos
{\bfseries\Large} % por defecto se usar�n caracteres de tama�o \Huge en negrita
{% contenido de la etiqueta
\color{blue! 30!black}
\Large\chaptertitlename\ %"Cap�tulo" o "Ap�ndice" en tama�o \Large en lugar de \Huge
\Large\thechapter .} % n�mero de cap�tulo en tama�o \Large
{2mm} % espacio m�nimo entre etiqueta y cuerpo
{\filright\color{blue! 30!black}} % texto del cuerpo alineado a la derecha
[\vspace{0.5mm}] % despu�s del cuerpo, dejar espacio vertical y trazar l�nea horizontal gruesa

%Secciones%
\titleformat{\section}[hang]
{\scshape\bfseries\Large}
{\color{blue! 30!black}\thesection. }
{0mm}{\filright\color{blue! 30!black}}[\vspace{0.5mm}]

%SubSecciones%
\titleformat{\subsection}[hang]
{\scshape\bfseries\large}
{\color{blue! 30!black}\thesubsection .}
{0mm}{\filright\color{blue! 30!black}}[\vspace{0.5mm}]

%SubsubSecciones%
\titleformat{\subsubsection}[hang]
{\scshape\bfseries}
{\color{blue! 30!black}\thesubsubsection .}
{0mm}{\filright\color{blue! 30!black}}[\vspace{0.5mm}]
%Fin de comandos%
%primera l�nea de cada p�rrafo%%%
\author{Hola}
\title{Esto es un ejemplo de \LaTeX en acci�n}
\date{La fecha de hoy}
\begin{document}%%%Aqu� empieza el documento%%%
\maketitle
\tableofcontents
\listoffigures
\part{Primer cap�tulo}
\AddToShipoutPicture{\BackgroundPic} 
\chapter{Primer cap�tulo}
{\quotation{
 Esto es una cita?}}
\section{Primera secci�n}
\subsection{Primera subsecci�n}
\subsubsection{Primera subsubsecci�n}
\lettrine{L}orem ipsum dolor sit amet, consectetur adipiscing elit. Duis tempus mattis felis, sed pulvinar leo sagittis at. Nulla ullamcorper aliquam eleifend. Nunc porta vulputate commodo. Nullam eu sem nec urna porttitor faucibus. Quisque vel feugiat mi. Ut orci neque, adipiscing eu fringilla facilisis, dictum et justo. Nam imperdiet congue metus at congue. Proin a ullamcorper augue. Vivamus a augue vel diam posuere sollicitudin in ut leo. Sed lectus neque, faucibus et pulvinar id, rutrum id ante. Donec quis tortor a tellus consectetur ultricies ut a mauris. 
%Texto de la secci�n
\begin {eqnarray}
 \label{eq1} %%%Comienzo
%de la ecuaci�n%%%
f:A \times M \rightarrow M \\
(\lambda, x) \rightarrow \lambda x \nonumber
\end {eqnarray} %%%Fin de la ecuaci�n%%%
\begin{enumerate}
\item Item 1
\item Item 2
\begin{enumerate}
\item Item 2.1
\item Item 2.2
\begin{enumerate}
\item Item 2.2.1
\item Item 2.2.2
\item Item 2.2.3
\end{enumerate}
\end{enumerate}
\item Item 3
\item Item 4
\end{enumerate}

\begin{itemize}
  \item Item 1
\item Item 2
\begin{itemize}
\item Item 2.1
\item Item 2.2
\begin{itemize}
\item Item 2.2.1
\item Item 2.2.2
\item Item 2.2.3
\end{itemize}
\end{itemize}
\item Item 3
\item Item 4
\end{itemize}


\begin{verbatim}
/**
	 * \brief Instancia un nuevo usuario.
	 *
	 * @param id the id
	 * @param nombre el nombre de pila del usuario.
	 * @param apellidos el apellido del usuario.
	 * @param string 
	 * @param rol el rol de usuario
	 * @param mail the mail
	 */
	public Usuario(Long id, String nombreUsuario, String nombre,
								 String apellidos, Rol rol, String mail) {

		this.nombre = nombre;
		this.apellidos = apellidos;
		this.rol = rol;
		this.email=mail;
		this.id=id;
		this.nombreUsuario = nombreUsuario;
	}
\end{verbatim}

\verb|this.nombre = nombre;|

\begin{verbatim*}
/**
	 * \brief Instancia un nuevo usuario.
	 *
	 * @param id the id
	 * @param nombre el nombre de pila del usuario.
	 * @param apellidos el apellido del usuario.
	 * @param string 
	 * @param rol el rol de usuario
	 * @param mail the mail
	 */
	public Usuario(Long id, String nombreUsuario, String nombre,
								 String apellidos, Rol rol, String mail) {

		this.nombre = nombre;
		this.apellidos = apellidos;
		this.rol = rol;
		this.email=mail;
		this.id=id;
		this.nombreUsuario = nombreUsuario;
	}
\end{verbatim*}

\imgIncl{13}{10}{img/060911scrumtoon.eps}
\imgIncl{13}{350}{img/060911scrumtoon.eps}



\end {document}